\section{Implementação}

\subsection{Organização do código}
A implementação foi estruturada em diferentes classes, de modo a separar as responsabilidades da aplicação. Essa organização permite que a representação dos dados, os cálculos matemáticos, a interface gráfica e a manipulação do display sejam tratados de forma independente.

A classe \texttt{Ponto} foi utilizada para representar coordenadas bidimensionais. Ela possui os atributos \texttt{x} e \texttt{y}, além de métodos de acesso e alteração desses valores. Dessa forma, as coordenadas utilizadas durante as transformações podem ser encapsuladas em objetos.

A classe \texttt{JanelaMundo} concentra os valores referentes aos limites do espaço de mundo. Seu construtor e seus métodos de alteração realizam verificações para garantir que os valores mínimos sejam menores que os valores máximos.

A classe \texttt{Transformacao} concentra os métodos responsáveis pelos cálculos matemáticos. Os métodos recebem os pontos e os parâmetros necessários para cada transformação e retornam novos objetos \texttt{Ponto} contendo os valores calculados.

A classe \texttt{Display} é responsável pela representação do display por meio de um \texttt{BufferedImage}. Já a classe \texttt{JanelaPrincipal} realiza a integração entre os componentes, controlando a interface, a entrada do mouse e a execução das transformações.

Por fim, a classe \texttt{Main} contém o ponto de inicialização da aplicação e cria a janela principal.

\subsection{Implementação das funções de transformação}
As funções de transformação foram implementadas de forma independente na classe \texttt{Transformacao}. Essa abordagem permite que cada operação matemática seja executada por um método específico.

Para a conversão da entrada para NDC \mbox{[0,1]}, foi implementado o método \texttt{inpToNdc01()}. Para NDC \mbox{[-1,1]}, foi utilizado o método \texttt{inpToNdc11()}.

As transformações entre NDC e a janela de mundo foram separadas nos métodos \texttt{ndc01ToUser()}, \texttt{userToNdc01()}, \texttt{ndc11ToUser()} e \texttt{userToNdc11()}. Também foram implementados os métodos \texttt{ndc01ToDc()} e \texttt{ndc11ToDc()}, responsáveis pela conversão das coordenadas normalizadas para as coordenadas do dispositivo.

Além dos cálculos, a classe possui o método privado \texttt{validarDimensoes()}, utilizado para verificar se as dimensões fornecidas para o display são válidas antes da realização das operações que dependem da largura e da altura.

\subsection{Implementação da entrada do mouse}
O rastreamento do mouse foi implementado na classe \texttt{JanelaPrincipal}, no método \texttt{configurarMouse()}, por meio de um \texttt{MouseInputAdapter} que sobrescreve três métodos: \texttt{mousePressed()}, chamado quando um botão é pressionado; \texttt{mouseMoved()}, chamado a cada movimento do mouse sobre o display; e \texttt{mouseDragged()}, chamado a cada movimento com o botão pressionado. O mesmo adaptador é registrado no display duas vezes, com \texttt{addMouseListener()} e com \texttt{addMouseMotionListener()}. Esse duplo registro é necessário porque o Swing entrega \texttt{mousePressed()} apenas aos \texttt{MouseListener}, enquanto \texttt{mouseMoved()} e \texttt{mouseDragged()} são entregues apenas aos \texttt{MouseMotionListener}.

Os três métodos encaminham o evento ao método \texttt{rastrearMouse()}, que obtém as coordenadas em relação ao display e verifica, por meio de \texttt{dentroDoDisplay()}, se elas estão dentro dos limites. Em caso positivo, os valores são encaminhados ao método \texttt{ativarPixelPorClique()}, que cria um objeto do tipo \texttt{Ponto} e o repassa aos métodos da classe \texttt{Transformacao}. Apesar do nome, mantido da versão inicial em que apenas o clique era tratado, esse método é executado para qualquer um dos três eventos.

Em \texttt{ativarPixelPorClique()}, o ponto de entrada é convertido para os dois sistemas NDC (\texttt{inpToNdc01()} e \texttt{inpToNdc11()}). Em seguida, o ponto do mundo é obtido a partir do NDC selecionado (\texttt{ndc01ToUser()} ou \texttt{ndc11ToUser()}) e convertido de volta para coordenadas de dispositivo pelo método auxiliar \texttt{mundoParaDispositivo()}. A atualização do pixel ativo e das informações da interface ocorre ao final desse método, de modo que cada nova posição do mouse produz uma nova ativação.

\subsection{Implementação da seleção do sistema NDC}
O sistema NDC em uso é armazenado no atributo \texttt{ndcAtual} da \texttt{JanelaPrincipal}, do tipo \texttt{TipoNDC}. O seletor é um \texttt{JComboBox} preenchido com \texttt{TipoNDC.values()}; seu \textit{listener} apenas atualiza \texttt{ndcAtual} com o item escolhido.

O método \texttt{mundoParaDispositivo()} centraliza o uso desse atributo: quando \texttt{ndcAtual} vale \texttt{ZERO\_A\_UM}, utiliza \texttt{userToNdc01()} seguido de \texttt{ndc01ToDc()}; caso contrário, utiliza \texttt{userToNdc11()} seguido de \texttt{ndc11ToDc()}. Esse método é compartilhado pelas duas formas de entrada, o que garante que o mouse e a consulta de ponto empreguem sempre o mesmo pipeline.

\subsection{Implementação da consulta de ponto do mundo}
O painel \textit{Consultar Ponto do Mundo} possui dois campos de texto, para $x_w$ e $y_w$, e o botão \textit{Mostrar no display}. Ao ser pressionado, o botão executa o método \texttt{consultarPontoMundo()}, que converte os valores digitados para \texttt{double}, cria o \texttt{Ponto} correspondente e obtém suas coordenadas de dispositivo por meio de \texttt{mundoParaDispositivo()}. O resultado é arredondado com \texttt{Math.round()} e, se estiver dentro do display, enviado a \texttt{drawPixel()}; caso contrário, o pixel não é ativado e um alerta é exibido (Subseção~\ref{subsec:impl-fora-limites}). Valores não numéricos são tratados por uma \texttt{NumberFormatException}, que resulta em uma mensagem de erro.

Os valores exibidos na parte inferior da janela são atualizados pelo método \texttt{atualizarLabels()}. Na consulta de ponto, esse método recalcula o ponto nos dois sistemas NDC (\texttt{userToNdc01()} e \texttt{userToNdc11()}) apenas para fins de exibição. Ao ser chamado, o método também apaga a mensagem de alerta que estiver sendo exibida, de modo que o alerta desaparece assim que uma coordenada válida volta a ser informada.

\subsection{Implementação do tratamento de coordenadas fora dos limites}
\label{subsec:impl-fora-limites}
A verificação dos limites é feita pelo método \texttt{dentroDoDisplay()}, que retorna verdadeiro apenas quando $0 \le x < 800$ e $0 \le y < 600$. Ela é aplicada em dois pontos: em \texttt{rastrearMouse()}, sobre a posição do mouse, e em \texttt{consultarPontoMundo()}, sobre o pixel obtido para o ponto do mundo, já arredondado.

Quando a coordenada está fora dos limites, é chamado o método \texttt{mostrarForaDosLimites()}, que limpa o display, substitui os textos de dispositivo, NDC \mbox{[0,1]}, NDC \mbox{[-1,1]} e mundo por ``fora dos limites'' e escreve a mensagem recebida em um \texttt{JLabel} de alerta, exibido em vermelho abaixo das coordenadas. Para o mouse, a mensagem informa a posição e a resolução do display, por exemplo ``Coordenada ($-15$, 320) fora dos limites do display 800x600.''. Para a consulta de ponto, a mensagem informa o ponto digitado e os limites da janela de mundo e também é apresentada em uma caixa de diálogo de aviso (\texttt{JOptionPane.WARNING\_MESSAGE}).

\subsection{Implementação do display}
A representação do display foi implementada na classe \texttt{Display} utilizando a classe \texttt{BufferedImage}. A imagem possui exatamente $800 \times 600$ posições, correspondentes à resolução que foi definida para a aplicação.

O método \texttt{drawPixel()} recebe as coordenadas inteiras do pixel que deverá ser ativado. Antes de realizar a alteração, o método verifica se as coordenadas estão dentro dos limites da imagem. Quando a posição é válida, o método \texttt{setRGB()} altera diretamente o valor de cor daquela posição. A cor utilizada para o pixel ativo é obtida a partir de \texttt{Color.GREEN}.

A limpeza do display é realizada pelo método \texttt{limpar()}, que preenche a imagem com a cor de fundo antes da ativação de uma nova posição. A imagem armazenada no \texttt{BufferedImage} é posteriormente apresentada pelo método \texttt{paintComponent()}, responsável por desenhá-la no componente gráfico.

\subsection{Implementação da interface e configuração da janela de mundo}
A interface gráfica foi construída utilizando componentes do Swing. A classe \texttt{JanelaPrincipal} utiliza um \texttt{BorderLayout}: o display ocupa a região central, as informações das coordenadas e a linha de alerta ficam na região inferior e um painel lateral, à direita, agrupa os painéis \textit{Janela do Mundo}, \textit{Sistema NDC} e \textit{Consultar Ponto do Mundo}. A janela é dimensionada com \texttt{pack()}, centralizada na tela e não pode ser redimensionada, o que mantém o display com exatamente $800 \times 600$ pixels.

Os campos de texto são utilizados para receber os valores de $x_{wmin}$, $x_{wmax}$, $y_{wmin}$ e $y_{wmax}$. Quando o usuário solicita a aplicação dos novos valores, o método responsável pela atualização realiza a conversão dos dados para \texttt{double} e cria uma nova instância de \texttt{JanelaMundo}.

A implementação também trata possíveis entradas inválidas. Valores que não podem ser convertidos para números ou que não respeitam os limites necessários são tratados por exceções e apresentados ao usuário por meio de mensagens na interface.

\subsection{Integração da implementação}
A integração entre as classes ocorre principalmente na \texttt{JanelaPrincipal}. Quando a posição do mouse muda ou um ponto do mundo é consultado, a classe cria o ponto inicial e utiliza os métodos da \texttt{Transformacao}, de acordo com o sistema NDC selecionado, para obter os valores correspondentes aos diferentes sistemas de coordenadas. Depois das transformações necessárias, a coordenada final é arredondada com \texttt{Math.round()}, verificada em relação aos limites do display e, se válida, enviada para o método \texttt{drawPixel()} da classe \texttt{Display}.

A separação entre essas responsabilidades permite que os cálculos de transformação permaneçam independentes da implementação gráfica. Ao adotar esse modelo, a classe \texttt{Transformacao} não precisa conhecer os componentes da interface ou a forma como o pixel é desenhado, enquanto \texttt{Display} não precisa realizar os cálculos matemáticos das coordenadas.

Essa estrutura mantém a implementação organizada e permite relacionar diretamente cada parte do código com sua respectiva função dentro da aplicação.
